% Generated draft: requires mathematical and source review.
\documentclass[11pt]{article}
\usepackage[margin=25mm]{geometry}
\usepackage{fontspec}
\setmainfont{Latin Modern Roman}
\usepackage{amsmath,amssymb,mathtools,unicode-math}
\setmathfont{Latin Modern Math}
\usepackage{enumitem,needspace,xurl,hyperref}
\setlength{\parindent}{0pt}
\setlength{\parskip}{5pt}
\setlength{\emergencystretch}{3em}
\newcommand{\N}{\mathbb{N}}
\newcommand{\Z}{\mathbb{Z}}
\newcommand{\Q}{\mathbb{Q}}
\newcommand{\R}{\mathbb{R}}
\newcommand{\C}{\mathbb{C}}
\newcommand{\F}{\mathbb{F}}
\newcommand{\A}{\mathbb{A}}
\newcommand{\PP}{\mathbb{P}}
\newcommand{\E}{\mathbb{E}}
\newcommand{\eps}{\varepsilon}
\newcommand{\dd}{\,\mathrm{d}}
\begin{document}
\section*{UnsolvedMath ID: 30005581}
\textbf{Source problem code:} catalogue entry
\subsection{Definitions and mathematical statement}
Let \(X\) be a compact smooth manifold with boundary. Its double is the closed smooth manifold
\[
D(X)=\operatorname{Dbl}(X,\partial X):=X\cup_{\partial X}X,
\]
formed by gluing two copies along the identity on their boundaries, using collars to give the double its smooth structure. Write
\[
\mathcal R^+(X)=\{g\in\operatorname{Met}^{\infty}(X):\operatorname{Scal}_g(x)>0\text{ for every }x\in X\},
\]
where \(\operatorname{Met}^{\infty}(X)\) denotes the smooth Riemannian metrics on \(X\). For a specified boundary condition \(\mathcal B\), put
\[
\mathcal R^+(X;\mathcal B)=\{g\in\mathcal R^+(X):g\text{ satisfies }\mathcal B\}.
\]
The usual smooth topology on metrics would give these sets their subspace topologies, but the supplied record does not explicitly specify the topology intended in its homotopy question.

A product boundary condition ordinarily means that there are \(\varepsilon>0\), a collar \(c:\partial X\times[0,\varepsilon)\longrightarrow X\), and a Riemannian metric \(h\) on \(\partial X\) such that
\[
c^*g=h+\dd t^2.
\]
This describes the interpretation with freely varying boundary metric and collar. The supplied excerpt does not state whether its boundary condition fixes either of these data. A mean-convex boundary condition is expressed using the mean curvature of \(\partial X\); its precise inequality depends on the normal and sign conventions, and on whether strict or weak mean convexity is required. Those choices are absent from the supplied record, so no particular inequality is imposed here.

\textbf{First source question.} Is the Doubling Conjecture true in its full intended generality? The supplied record names this conjecture but does not state it. In particular, it does not specify the implication or equivalence between positive scalar curvature on \(D(X)\) and metrics on \(X\), its boundary hypotheses, or the class of manifolds quantified over. Consequently, a precise quantified version of this question cannot be recovered from the supplied evidence alone.

\textbf{Second source question.} For the compact manifolds \(X\) allowed by the source, determine the homotopy types of
\[
\mathcal R^+(X;\mathcal B_1)
\quad\text{and}\quad
\mathcal R^+(X;\mathcal B_2),
\]
where \(\mathcal B_1\) and \(\mathcal B_2\) are the source's boundary conditions numbered (1) and (2). The curated statement describes these as product-type and corresponding mean-convex boundary conditions. Their exact definitions, the permitted dimensions and other manifold hypotheses, and any fixed boundary data must be recovered before this becomes a fully specified homotopy-type problem.

\textbf{Contextual qualifications.} The original excerpt states that, in an unspecified case,
\[
M=\operatorname{Dbl}(X,\partial X)=Y\times S^1
\]
admits a positive-scalar-curvature metric, and that the source's Theorem 1 concerning obstructions to product-type boundary conditions fails in that case. The excerpt does not define \(Y\), describe that case, or supply Theorem 1; this assertion is therefore retained only as source-specific context, not as a claim for arbitrary \(Y\). It also reports proofs in many cases using surgery techniques from its reference [3], without supplying that reference or making the antecedent of ``this'' precise.
\subsection{Short English statement}
The source asks whether the Doubling Conjecture holds generally and asks for the homotopy types of spaces of positive-scalar-curvature metrics with product-type or mean-convex boundary conditions. The conjecture itself and the exact boundary conditions are missing from the supplied excerpt, so the formulation remains incomplete. The supplied open status and reported special cases require independent review.
\subsection{Sources}
\begin{itemize}
\item Local UnsolvedMath registry, permanent ID \(30005581\), exact source problem code \texttt{catalogue entry}, title ``Doubling and Homotopy of Positive Scalar Curvature Metrics.'' The record belongs to the Oberwolfach Reports Open Problems collection and supplies proposed year 2023; no proposer is named.
\item Supplied report citation: \emph{Differentialgeometrie im Grossen (2024)}, Oberwolfach Reports, DOI 10.4171/owr/2023/29, pp.\ 1617--1670. \url{https://doi.org/10.4171/owr/2023/29}. This is also the supplied literature source labelled ``original OWR report.'' The report contents have not been independently consulted for this draft.
\item Supplied literature source labelled ``Rosenberg--Weinberger'': \url{https://arxiv.org/abs/2201.01263}. No paper title or precise theorem statement is supplied.
\item Supplied literature source labelled ``Frenck'': \url{https://arxiv.org/abs/2604.12901}. No paper title or precise theorem statement is supplied.
\item Literature assessment dated exactly \texttt{2026-08-22}: the supplied curation labels the problem open and reports that the doubling conjecture is proved in dimension three and under a split fundamental-group condition, while its general form and the accompanying homotopy-type question remain unresolved. This is supplied dated curation, not an independently verified literature assessment; the precise hypotheses of the reported special cases are absent.
\item The record labels its statement \texttt{corrected\_verified} and reports that the official report was checked directly. That verification claim belongs to the supplied record and is not an independent verification by this draft. The original excerpt also cites a Theorem 1 and a reference [3], neither of which is supplied in identifiable form.
\end{itemize}
\end{document}
